\section*{Resumo}\markboth{Resumo}{}
Este trabalho avalia a contribuição preditiva incremental de representações de letras para estimar a popularidade observada no Spotify, em conjunto com atributos acústicos e metadados. O ponto de partida é um catálogo público com 1.159.764 faixas e 82 gêneros. Desse catálogo foi selecionada uma amostra probabilística estratificada de 200.000 faixas, incluindo integralmente as 499 com popularidade de pelo menos 80. A recuperação local de letras no LRCLIB resultou em 77.535 faixas aceitas. A comparação descritiva revelou disponibilidade textual seletiva: a popularidade média do catálogo é 18,38, enquanto a média das aceitas é 23,98. A modelagem foi, portanto, delimitada ao domínio analítico observado, sem pressupor representatividade do catálogo. Compararam-se quatro configurações de regressão Ridge: atributos acústicos e metadados; esses atributos acrescidos de 21 descritores lexicais e estruturais; acrescidos de embeddings multilingues; e acrescidos de ambas as representações textuais. A avaliação utilizou validação cruzada agrupada aninhada, com cinco partições externas e três internas, preservando 21.893 componentes operacionais de artista, composição, versão e texto. O erro médio absoluto passou de 8,3525 no Baseline para 8,1151 no Combinado, redução de 2,84\%. O acréscimo dos descritores aos embeddings produziu ganho menor, de 0,0134 ponto, inconclusivo na análise auxiliar de incerteza. As 128 faixas aceitas com popularidade de pelo menos 80 foram todas subestimadas. Os resultados indicam contribuição textual modesta e não uniforme na avaliação realizada; não demonstram causalidade, previsão prospectiva de sucesso ou generalização automática ao catálogo Spotify.

\vspace{.6cm}
{\raggedright\noindent\textbf{Palavras-chave:} popularidade musical; letras; amostragem estratificada; regressão Ridge; embeddings; validação cruzada agrupada.\par}
